\section{Proposed Technique}

\subsection{Approach Overview}
% [proposal] Describe different components and stages of your technique using text/figure/example/flowchart/etc. 
Our initial design was a test-level GNN-TCP system: construct a graph connecting changed methods to covered tests, propagate change impact through the graph, and rank tests by predicted failure probability. In practice, robust method-level dynamic coverage was not available for the full LRTS setting. The implemented technique is therefore a static build-failure baseline. It has three stages: (1) construct one static call graph per project, (2) encode each project graph using a two-layer Graph Convolutional Network (GCN), and (3) concatenate the resulting project embedding with per-build CI metadata to predict whether a build contains at least one failing test class.

\subsection{Static Graph Construction}

For each project $p$, we represent the source code as a static directed call graph
\[
G_p=(V_p,E_p),
\]
where each node $v\in V_p$ is a method and each directed edge $(u,v)\in E_p$ indicates that method $u$ invokes method $v$. The graph is extracted from Java Abstract Syntax Trees by identifying \texttt{MethodDeclaration} and \texttt{MethodInvocation} nodes.

Let $d^-_v$ and $d^+_v$ denote the in-degree and out-degree of node $v$ in the directed call graph. The implemented node feature vector is purely structural:
\[
x_v =
\begin{bmatrix}
\log(1+d^-_v) &
\log(1+d^+_v) &
\log(1+d^-_v+d^+_v) &
1
\end{bmatrix}.
\]
Stacking all node features gives $X_p\in\mathbb{R}^{|V_p|\times 4}$.

Although the extracted call graph is directed, the current implementation symmetrizes it before message passing. If $A_p$ is the directed adjacency matrix, the model uses
\[
\tilde{A}_p = A_p + A_p^\top + I,
\]
where $I$ adds self-loops. It then applies symmetric normalization
\[
\hat{A}_p = \tilde{D}_p^{-1/2}\tilde{A}_p\tilde{D}_p^{-1/2},
\qquad
\tilde{D}_{p,ii}=\sum_j \tilde{A}_{p,ij}.
\]
This normalization is mathematically standard for GCNs, but it also removes caller-callee directionality, which later becomes an important limitation.

\subsection{GCN Build Failure Model}

The graph encoder uses two graph convolution layers:
\[
H_p^{(1)}=\operatorname{ReLU}(\hat{A}_pX_pW_0+b_0),
\]
\[
H_p^{(2)}=\operatorname{ReLU}(\hat{A}_pH_p^{(1)}W_1+b_1).
\]
The project embedding is the mean of final node states:
\[
g_p=\frac{1}{|V_p|}\sum_{v\in V_p}H_{p,v}^{(2)}.
\]
For a build $i$ belonging to project $p(i)$, let $z_i$ be its normalized CI metadata vector. The classifier predicts a logit
\[
s_i = f_\theta([g_{p(i)};z_i]),
\]
where $[\cdot;\cdot]$ denotes concatenation and $f_\theta$ is a two-layer multilayer perceptron. The predicted failure probability is
\[
\hat{y}_i=\sigma(s_i)=\frac{1}{1+\exp(-s_i)}.
\]
The target label is
\[
y_i=\mathbb{1}[\texttt{num\_fail\_class}_i>0].
\]

Because positive and negative builds are imbalanced, training uses weighted binary cross-entropy:
\[
\mathcal{L}
=-\frac{1}{N}\sum_{i=1}^{N}
\left[
\alpha y_i\log \hat{y}_i
+(1-y_i)\log(1-\hat{y}_i)
\right],
\qquad
\alpha=\frac{N_-}{N_+},
\]
where $N_+$ and $N_-$ are the numbers of positive and negative training builds. A build is classified as failing when $\hat{y}_i\ge \tau$, where $\tau$ is selected on the validation split to maximize F1 over thresholds in $[0.2,0.8]$.


\subsection{Important Consequence of the Formulation}

The formulation above implies that all builds from the same project share the same graph embedding $g_p$. Thus, for two builds $i$ and $j$ from the same project,
\[
p(i)=p(j)\quad\Longrightarrow\quad g_{p(i)}=g_{p(j)}.
\]
Any build-specific variation must therefore come only from $z_i$, the CI metadata vector. The graph cannot encode which files changed, which methods were executed, or which tests covered the changed code for a particular build. This mathematical property is central to the empirical failure analyzed in Section~\ref{sec:evaluation}: the graph representation is static and project-level, so it can easily become a soft project identifier rather than a build-specific dependency signal.
